\section{Actual ideal module comparison under open restriction}
\begin{lemma}[Whole ideal membership under the actual open section isomorphism]
\label{lem:actual-ideal-open-section-membership}
\uses{lem:actual-ideal-module-whole-kernel}
\leanok
\lean{AlgebraicGeometry.Scheme.IdealSheafData.isLocallyMem_comap_iff_appIso_inv}
For an actual open immersion $g:X\to Y$, a section of the true geometric
ideal comap over any entire original open $W\subset X$ corresponds
exactly, under the genuine original structure-ring section isomorphism,
to an original ideal-module section over the entire image open $gW$.
\end{lemma}
\begin{proof}
For each actual affine $U\subset W$, the genuine open-comap affine ideal
formula is the inverse-image ideal under the original inverse section
isomorphism. Its naturality computes restriction as the original target
restriction. This proves the implication from original ideal membership
to comap membership. Conversely, for every original affine $V\subset gW$,
its true preimage is affine because $V$ lies in the open immersion's
range. Its image equals the whole original $V$, and its preimage lies
in $W$. Apply actual comap membership there and the genuine affine ideal
formula. Transport to the original affine $V$ through the true ideal
restriction map. Naturality of the actual section isomorphisms and the
whole restriction-composition law give precisely the original target
ideal membership. Thus the equivalence holds on arbitrary whole opens,
without flatness, Cartier or finiteness assumptions. Its module-sheaf
isomorphism and canonical-pullback gluing are subsequent constructions.
\end{proof}

\begin{lemma}[True whole ideal-module sheaf isomorphism under open restriction]
\label{lem:actual-ideal-open-restriction-module-iso}
\uses{lem:actual-ideal-open-section-membership}
\leanok
\lean{AlgebraicGeometry.Scheme.IdealSheafData.openRestrictionSectionsEquiv,AlgebraicGeometry.Scheme.IdealSheafData.openRestrictionModuleIso,AlgebraicGeometry.Scheme.IdealSheafData.openRestrictionModuleIso_app}
The actual original ideal module restricted along an arbitrary open
immersion is isomorphic to the actual module sheaf of its entire
geometric ideal comap. On every whole original open its underlying
included section map is precisely the original scheme section isomorphism.
\end{lemma}
\begin{proof}
The original section isomorphism and its true inverse restrict to the
actual original ideal sections by the preceding whole membership
equivalence. Their genuine scalar actions are intertwined: the restricted
module's action uses the original inverse ring comparison, so original
multiplicativity and the true inverse identities give linearity over
the original source section ring. Naturality is exactly the original
section-isomorphism naturality on every original inclusion. Thus the
section linear equivalences define an actual presheaf-module isomorphism;
the genuine fully faithful inclusion of actual sheaves of modules lifts
it to the original whole sheaf isomorphism. The explicit full section
formula follows from this actual construction. No quasicoherence,
Cartier, flatness or comparison premise is used.
\end{proof}


\begin{lemma}[Canonical whole ideal-module pullback isomorphism for an open immersion]
\label{lem:actual-ideal-open-canonical-pullback-iso}
\uses{lem:actual-ideal-open-restriction-module-iso,lem:actual-ideal-pullback-unit-sections}
\leanok
\lean{AlgebraicGeometry.Scheme.IdealSheafData.openRestrictionModuleIso_restrictAdjunction,AlgebraicGeometry.Scheme.IdealSheafData.openRestrictionModuleIso_comp_toModulesComapPullbackMap,AlgebraicGeometry.Scheme.IdealSheafData.toModulesComapPullbackMap_isIso_of_isOpenImmersion,AlgebraicGeometry.Scheme.IdealSheafData.openPullbackModuleIso}
For every original ideal and arbitrary actual open immersion, its genuine
canonical module pullback comparison with the module of its geometric
ideal comap is an isomorphism. Under the true restriction--pullback
comparison it is the whole open-restriction isomorphism just constructed.
\end{lemma}
\begin{proof}
Apply the actual restriction--pushforward adjunction to the constructed
whole restriction isomorphism. Its adjoint on any original open first
restricts the original ideal section to the whole image of its preimage,
then applies the genuine original open section isomorphism. The actual
scheme identity for these ring maps says that this composite is exactly
the original whole scheme section map. Thus the adjoint is the canonical
ideal-module pushforward map. The true uniqueness isomorphism between
the restriction and pullback left adjoints identifies their units.
Naturality of the original adjunction and the proved canonical
pullback-unit identity therefore identifies the restriction isomorphism
with that uniqueness isomorphism followed by the canonical pullback map.
Injectivity of the actual adjunction hom equivalence proves the whole
sheaf equality. Both the restriction isomorphism and the uniqueness
comparison are isomorphisms, hence so is the actual canonical map.
Taking its genuine inverse gives the stated canonical sheaf isomorphism.
No Cartier, quasicoherence, flatness or finiteness premise occurs.
\end{proof}
